\documentclass[12pt]{article}
\usepackage[margin=0.3in,top=0.25in,bottom=0.4in]{geometry}
\makeatletter
\renewcommand{\@maketitle}{\newpage\null\vskip -3.5em\begin{center}{\LARGE \@title \par}\end{center}\par\vskip 0.2em}
\makeatother
\usepackage{amsmath,amssymb}
\usepackage{booktabs}
\usepackage{hyperref}
\setlength{\parskip}{0.6em}
\setlength{\parindent}{0pt}
\newcommand{\xpo}{x_{\mathbb{P}}}

\title{Two fixes in the diffractive (TMD) integrand}
\author{}
\date{}

\begin{document}
\maketitle

\section*{In short}

Compared with the derivation in Ref.~[3] (Caucal, Gimeno-Estivill, Iancu, Lappi, Salazar,
arXiv:2606.04169), the diffractive (TMD) part of the code has two errors:
\begin{enumerate}
\item[(a)] it uses the \textbf{quark} dipole amplitude $\mathcal{N}$ instead of the \textbf{gluon}
  dipole amplitude $\mathcal{N}_g=2\mathcal{N}-\mathcal{N}^2$;
\item[(b)] it integrates over $\mathrm{d}\xpo$ instead of $\mathrm{d}\ln\xpo$, i.e.\ with an extra factor
  $\xpo$ in the integrand.
\end{enumerate}
Both make the diffractive TMD cross section too small. The exclusive part is correct. After the fix,
all diffractive results have to be recomputed.

\section*{1. The formula of Ref.~[3]}

At fixed photon energy, Eqs.~(3.38) and (3.43)--(3.45) of Ref.~[3] give
\[
  \frac{\mathrm{d}\sigma}{\mathrm{d}z_1\mathrm{d}z_2\,\mathrm{d}^2K\,\mathrm{d}^2P\,\mathrm{d}\ln(1/x)}
  =\mathcal{H}\,\frac{\mathrm{d}\,xG_{\mathbb{P}}}{\mathrm{d}^2K},\qquad
  \frac{\mathrm{d}\,xG_{\mathbb{P}}}{\mathrm{d}^2K}=\frac{S_\perp(N_c^2-1)}{8\pi^4(1-x)}\,G^2,
\]
\[
  G=\Omega^2\int_0^\infty\mathrm{d}R_\perp\,R_\perp\,J_2(K_\perp R_\perp)\,K_2(\Omega R_\perp)\,
  \mathcal{N}_g(R_\perp,\xpo),\qquad \Omega^2=\frac{x}{1-x}K_\perp^2 .
\]
The code agrees with this in the hard factor $\mathcal{H}$, in $\Omega^2$, in the radial integral, in
$1/(1-x)$ and in the prefactor $S_\perp(N_c^2-1)/(8\pi^4)$. It differs in two places.

\section*{2. (a) Quark vs.\ gluon dipole}

In the diffractive process the $q\bar q g$ system scatters effectively as a gluon-gluon dipole, so $G$
contains $\mathcal{N}_g=2\mathcal{N}-\mathcal{N}^2$ (large $N_c$). The code called
\texttt{dipole->N(r, x)}, the quark dipole. AmplitudeLib already provides the gluon dipole as
\texttt{N\_A(r, x)}.

\textbf{Why $\mathcal{N}_g>\mathcal{N}$.} A gluon carries more colour charge than a quark, so a
gluon--gluon dipole interacts more strongly with the target than a quark--antiquark dipole of the same
size. At large $N_c$ a gluon behaves like a quark--antiquark pair at the same point, so the gluon
dipole passes through the target without scattering only if both of its quark dipoles do. The
``survival'' amplitudes $S=1-\mathcal{N}$ therefore multiply:
\[
  S_g=S^2\quad\Longrightarrow\quad 1-\mathcal{N}_g=(1-\mathcal{N})^2\quad\Longrightarrow\quad
  \mathcal{N}_g=2\mathcal{N}-\mathcal{N}^2\,.
\]
The difference is $\mathcal{N}_g-\mathcal{N}=\mathcal{N}(1-\mathcal{N})\geq0$, because
$0\leq\mathcal{N}\leq1$; the two are equal only at $\mathcal{N}=0$ (no scattering) and
$\mathcal{N}=1$ (saturation).
\begin{itemize}
\item \textbf{Small dipoles} (dilute, $\mathcal{N}\ll1$): $\mathcal{N}_g\approx2\mathcal{N}$. The factor 2
  is the ratio of colour charges, $C_A/C_F=9/4$, which is 2 at large $N_c$. In the cross section
  ($\propto G^2\propto\mathcal{N}_g^2$) this is a factor $\approx4$.
\item \textbf{Large dipoles} (saturated, $\mathcal{N}\to1$): both go to 1, since nothing can scatter with
  more than 100\% probability, and the gain disappears.
\end{itemize}
Equivalently, the gluon dipole saturates earlier: $Q_{s,g}^2\approx(C_A/C_F)\,Q_s^2\approx2\,Q_s^2$.

The fix therefore always increases the cross section, by a factor between 1 and 4. Computed with our dipole amplitudes, the factor $(G[\mathcal{N}_g]/G[\mathcal{N}])^2$ is:
\begin{center}
\begin{tabular}{cccc}
\toprule
$K_\perp$ [GeV] & proton, $\xpo=10^{-2}$ & proton, $\xpo=10^{-3}$ & Pb centre, $\xpo=10^{-2}$ \\
\midrule
0.3 & 1.2--2.0 & 1.25--1.95 & 1.06--1.4 \\
1.0 & 2.8--3.5 & 2.6--3.3   & 1.7--2.7 \\
3.0 & 3.8--4.0 & 3.7--3.9   & 3.4--3.8 \\
\bottomrule
\end{tabular}
\end{center}
(ranges over $x=0.1$--$0.9$). It approaches 4 for small dipoles ($\mathcal{N}_g\approx2\mathcal{N}$) and
1 near saturation. It is smaller for Pb than for the proton, so the diffractive $R_{pA}$ changes as
well: nuclear suppression becomes stronger.

\section*{3. (b) The $\xpo$ measure}

For given photon momentum and quark kinematics, $x=x_{q\bar q}/\xpo$ with
$x_{q\bar q}=M_{q\bar q}^2/W^2$ fixed. Hence $\ln(1/x)=\ln\xpo-\ln x_{q\bar q}$ and
\[
  \mathrm{d}\ln(1/x)=\mathrm{d}\ln\xpo\quad\Longrightarrow\quad
  \sigma=\int\mathrm{d}\ln\xpo\;f=\int\frac{\mathrm{d}\xpo}{\xpo}\;f\,,\qquad
  f=\mathcal{H}\,\frac{\mathrm{d}\,xG_{\mathbb{P}}}{\mathrm{d}^2K}=\frac{\mathrm{d}\sigma}{\mathrm{d}\ln\xpo}\,.
\]
The code samples $\xpo$ logarithmically, $\xpo=x_{\min}e^{uL}$ with $u\in[0,1]$ and
$L=\ln(0.1/x_{\min})$. The Jacobian for $\int\mathrm{d}\ln\xpo$ is $L$, while the code used
$\xpo L$ (\texttt{x\_po * log\_xpo\_range}), i.e.\ it computed $\int\mathrm{d}\xpo\,f$: an extra factor
$\xpo\leq0.1$ in the integrand. Correct: \texttt{log\_xpo\_range} only.

The same factor is in \texttt{integrand\_diffractive\_fixed\_qp} (charm at fixed $q^+$).

\section*{4. The $\xpo$ figure}

The figure shows $\mathrm{d}\sigma/\mathrm{d}\ln\xpo$ against $\ln\xpo$, so that the area under the
curve is the cross section, $\sigma=\int\mathrm{d}\ln\xpo\,\mathrm{d}\sigma/\mathrm{d}\ln\xpo$.

\begin{itemize}
\item \textbf{So far:} the fixed-$\xpo$ program returns $f$, which by Ref.~[3] already is
  $\mathrm{d}\sigma/\mathrm{d}\ln\xpo$. The plotting scripts (\texttt{xpom.py},
  \texttt{charm\_fixed\_qp\_xpom.py}) multiplied it by $\xpo$, treating it as
  $\mathrm{d}\sigma/\mathrm{d}\xpo$. The area under the figure did equal the cross section the code
  computed, $\int\mathrm{d}\ln\xpo\,(\xpo f)$: figure and code were consistent with each other, but both
  carried the extra $\xpo$. This is why the curves rise with $\xpo$ towards the cut at $0.1$.
\item \textbf{After the fix:} the integrated cross section is $\int\mathrm{d}\ln\xpo\,f$, and the
  scripts plot $f$ without the extra $\xpo$. The area is again the cross section, now the correct one.
  The curves become much flatter, so more of the cross section comes from small $\xpo$. The earlier
  statement that 86--98\% of the diffractive TMD part comes from $\xpo>0.01$ is an artefact of the
  extra $\xpo$ and should not be used.
\end{itemize}

\section*{5. Effect on the results}

\begin{itemize}
\item \textbf{Diffractive (TMD) part:} larger, by a factor 1.2--4 from (a) and by a larger,
  $\xpo$-dependent factor from (b); the exact size only comes out of the reruns.
\item \textbf{Exclusive part:} unchanged (it correctly uses the quark dipole and has no $\xpo$
  integral).
\item \textbf{$R_{pA}$ of the diffractive part:} changes, because (a) affects Pb and the proton
  differently.
\item \textbf{Manuscript:} Eqs.~(12)--(15) should follow Ref.~[3]: $\Omega^2$ (not $\Omega$), the radial
  integral $\int\mathrm{d}R\,R$ (not $\int\mathrm{d}^2R$), $1/(1-x)$ (not $1/(x(1-x))$) and
  $\mathcal{N}_g$.
\end{itemize}

\section*{6. Plan}

\begin{enumerate}
\item In \texttt{src/integrands.cpp}: \texttt{N} $\to$ \texttt{N\_A} in the three diffractive
  integrands (\texttt{integrand\_diffractive}, \texttt{integrand\_diffractive\_xpom} and the fixed-$q^+$
  kernel), and \texttt{x\_po * log\_xpo\_range} $\to$ \texttt{log\_xpo\_range} in the two $\xpo$
  integrals.
\item In \texttt{xpom.py} and \texttt{charm\_fixed\_qp\_xpom.py}: no extra factor $\xpo$; a header
  line in new files so that old and new results are never mixed.
\item Rerun all diffractive results: the fixed-$q^+$ charm results locally (about an hour), the
  $D^0$ results on Roihu (central values, scale variations, p--Pb, STARlight, $\xpo$ and the HymnD and
  BK bands), together with the $\sqrt2$ flux and $\alpha_s$ fixes in one rebuild there.
\item Regenerate all plots, and the diffractive input of the inclusive subtraction figure.
\end{enumerate}

\end{document}
